\documentclass[10pt]{article}
\usepackage{icmlko}

\icmlmeta{IP 파운데이션 월드모델}{245}{흉내 그 이상은 없다}{안쪽을 유보에 가깝게 만들려는 손질 여섯이 전부 실패했다 --- 제일 좋은 안쪽은 배포를 그대로 베낀 것이다}

\begin{document}

\twocolumn[
\icmlheader

\begin{abstract}
노트 244는 안쪽 눈금이 경계 바로 앞 \textbf{한 조각}일 때 제일 잘 읽는다는 것을
쟀고($|\Delta| \le 0.010$ 칸에서 유보와 spearman $+$0.51), 그 편향의 정체를
숙제로 남겼다. 세 가설을 세워 다 재고 다 반증했다. \textbf{①~라벨 누적}
(노트 209) --- 학습 구간에서 시작연도와 라벨의 상관이 0.30 넘는 도메인은
게임 · 아이돌 둘뿐이고, 그 둘을 빼면 오히려 $+$0.33 으로 \emph{내려간다}.
\textbf{②~도메인 섞임} --- 조각마다 섞임이 다르니 무게를 유보 섞임으로
고정했다. $+$0.47 로 내려간다. \textbf{③~안 보이는 도메인} --- 도서 · 아이돌 ·
팝업은 안쪽 조각에 \textbf{한 건도 없는데} 유보의 9.7\%이고, 만화는 그 반대다
(안쪽 6.3\%, 유보 0\%). 유보를 안쪽이 볼 수 있는 여섯으로 좁혀 견줬다.
$+$0.35 로 내려간다. \textbf{여섯 손질이 전부 흉내보다 못하다.} 남는 결론은
소극적이지만 분명하다 --- \textbf{제일 좋은 안쪽은 배포를 그대로 베낀
것이고}(경계 전부로 배워 경계 직후를 맞힌다), 거기서 더 영리해지려는 시도는
전부 손해다. 안쪽 신뢰도 상한 0.34 는 고칠 수 있는 결함이 아니라 \textbf{이
자료에서의 천장}으로 본다. 그러니 $|\Delta| \le 0.010$ 자리는 안쪽으로
정하지 않는다 --- 유보로도 정하면 안 되니(그건 엿보기다), \textbf{그 자리는
다투지 않는 것}이 답이다.
\end{abstract}
]

\section{숙제 하나, 가설 셋}

노트 244가 남긴 첫 줄은 ``왜 2022가 다른가 --- 시대 편향의 정체''였다.
조각을 늘리면 나빠지고 줄이면 좋아진다는 것은 편향이라는 뜻인데, 무엇의
편향인지는 안 밝혔다. 셋을 세웠다.

\section{① 라벨 누적이 아니다}

노트 209는 라벨이 누적인 도메인이 여섯이고 \texttt{-시작일} 하나만으로도
판이 0.2513 이라고 적었다. 누적이면 2022 레코드는 2026까지 4년을 쌓았고
2024 레코드는 2년만 쌓았으니, 옛 조각일수록 라벨이 \emph{나이}로 설명되고
축은 덜 중요해진다 --- 그럴듯한 편향이다.

먼저 지금 판에서 다시 쟀다.

\begin{center}\small
\begin{tabular}{lr@{\qquad}lr}
\hline
도메인 & $\rho$ & 도메인 & $\rho$ \\
\hline
아이돌 & $-$0.405 & 만화 & $+$0.159 \\
게임 & $-$0.363 & 세계애니 & $-$0.135 \\
펀딩 & $+$0.240 & 도서 & $+$0.132 \\
웹툰 & $+$0.234 & 모바일 & $-$0.126 \\
애니 & $+$0.175 & & \\
\hline
\end{tabular}
\end{center}

0.30을 넘는 것이 \textbf{둘뿐}이다(노트 209의 $-$0.648 · $-$0.579 보다 훨씬
약하다 --- 그때는 옛 라벨이었다). 그 둘을 빼고 나머지 일곱으로만 견주니
작은 칸 상관이 $+$0.511 에서 \textbf{$+$0.327}로 내려간다. 누적을 없앴는데
나빠졌다. 반증.

\section{② 도메인 섞임이 아니다}

조각마다 무엇이 들어 있는지 보니 눈에 띄는 것이 있었다.

\begin{center}\small
\begin{tabular}{lrrrr}
\hline
도메인 & 2022 & 2023 & 2024 & 유보 \\
\hline
웹툰 & 39.4\% & 39.0\% & 40.1\% & 27.7\% \\
애니 & 19.4\% & 13.8\% & 16.9\% & 23.6\% \\
세계애니 & 12.9\% & 18.1\% & 15.5\% & 11.7\% \\
모바일 & 12.3\% & 12.9\% & 15.3\% & 17.2\% \\
\textbf{만화} & 6.3\% & 5.1\% & 1.9\% & \textbf{0.0\%} \\
\textbf{도서} & \textbf{0.0\%} & \textbf{0.0\%} & \textbf{0.0\%} & 6.4\% \\
\textbf{팝업} & \textbf{0.0\%} & \textbf{0.0\%} & \textbf{0.0\%} & 2.3\% \\
\textbf{아이돌} & \textbf{0.0\%} & \textbf{0.0\%} & \textbf{0.0\%} & 1.0\% \\
\hline
\end{tabular}
\end{center}

섞임이 다르다. 무게를 유보 섞임으로 고정해 조각마다 같은 저울을 쓰게 했다.
작은 칸이 $+$0.357 $\to$ $+$0.269(조각 셋) · $+$0.511 $\to$
\textbf{$+$0.467}(한 조각)로 \emph{내려간다}. 반증.

\section{③ 안 보이는 도메인이 아니다}

위 표의 굵은 줄이 더 근본적인 문제를 말한다. \textbf{도서 · 팝업 · 아이돌은
안쪽 조각에 한 건도 없다}(2022--2024에 20건 이상인 적이 없다). 그런데 유보의
9.7\%다. 만화는 정반대로 안쪽에는 있고 유보에는 없다. \textbf{안쪽은 유보의
10분의 1을 원리적으로 못 본다.}

그렇다면 안쪽이 볼 수 있는 여섯(게임 · 모바일 · 세계애니 · 애니 · 웹툰 ·
펀딩)으로 유보를 좁혀 견주면 눈금이 맞아야 한다. 좁혔다. $+$0.511 $\to$
\textbf{$+$0.352}. 반증.

\section{여섯이 다 흉내보다 못하다}

\begin{figure}[t]
\centering
\includegraphics[width=\columnwidth]{figs/mimic.pdf}
\caption{안쪽을 손보는 여섯 가지. 점선이 아무것도 안 하고 배포를 그대로
베낀 것($+$0.511)이다.}
\end{figure}

\begin{center}\small
\begin{tabular}{lr}
\hline
안쪽 방식 & 작은 칸 $\rho$ \\
\hline
\textbf{2024 한 조각 (배포 흉내)} & $\mathbf{+0.511}$ \\
무게를 유보 섞임으로 고정 & $+$0.467 \\
해 셋 평균 & $+$0.357 \\
안쪽이 보는 여섯만 & $+$0.352 \\
누적 도메인 뺌 & $+$0.327 \\
반해 여섯 평균 & $+$0.214 \\
해 다섯 평균 & $+$0.203 \\
\hline
\end{tabular}
\end{center}

여섯 손질이 전부 아래에 있다. 이 정도로 방향이 한결같으면 개별 가설의
반증을 넘어서는 이야기다 --- \textbf{안쪽을 ``더 낫게'' 만들려는 조작이
곧 배포에서 멀어지는 조작}이었다.

\section{천장으로 받아들인다}

노트 244가 잰 안쪽 신뢰도 상한 0.34 를, 고칠 수 있는 결함이 아니라
\textbf{이 자료에서의 천장}으로 본다. 근거는 이번 여섯 번의 실패이고,
반증 가능한 형태로 적어 둔다 --- \emph{안쪽을 손봐서 0.34 를 넘기는 방법이
있다면 그것은 이 노트를 반증한다.}

따라오는 규약이 소극적이라 마음에 안 들지만 정직하다. $|\Delta| \le 0.010$
자리는 \textbf{안쪽으로 못 정하고}(신뢰도 0.34), \textbf{유보로도 정하면 안
된다}(그건 엿보기다). 그러면 남는 길은 하나다 --- \textbf{그 자리를 안
다툰다.} 노트 216 · 224 · 234 · 241에서 거부한 넷이 전부 그 자리였다.
거부가 옳았다는 것을 이제 눈금으로 안다.

\paragraph{남은 것}
\begin{center}\small
\begin{tabular}{ll}
\hline
갈래 & 왜 \\
\hline
0.010 위 & 거기는 안쪽이 읽는다(신뢰도 상한 0.16$\to$0.34, 큰 칸 부호 5/5) \\
도서 · 팝업 · 아이돌 & 안쪽에 한 건도 없다. 수집이 시대를 못 채웠다 \\
\texttt{entry\_friction} & 빼면 $-$0.0637 --- 판에서 제일 무거운 축 \\
만화 & 안쪽 6.3\%, 유보 0\% --- 유보에 없는 도메인을 왜 학습하나 \\
\hline
\end{tabular}
\end{center}

\paragraph{규약} 후보를 재기 전에 \textbf{크기부터 어림한다.} 0.010 아래로
예상되면 재지 않는다 --- 재도 정할 수 없기 때문이다. 0.010 위라면 경계 앞
한 조각으로 재고, 그 수는 근거가 된다.

\end{document}
